% ============================================================
%  06.tex --- Implementation: from Model to Feynman Rule
% ============================================================

% ==============================================================
\chapter{Implementation: From \texttt{Model} to Feynman Rule}
\label{ch:internals}
% ==============================================================

Chapter~\ref{ch:feynpy} served as a user guide, whilst Chapter~\ref{ch:cas} outlined the libraries on which FEYNPY is based. Now that we know how to use the toolkit and understand its language which, although it is Python, may not be familiar to everyone, we can move on to understanding how it works internally. How do we get from the declarative Python objects that the user writes to the Symbolica expression that emerges as a Feynman rule?

We physicists, in understanding the origin and physical meaning of a theory, make use of many conventions and simplifications (which, moreover, differ from physicist to physicist), and the code must recognise these conventions. When considering the Standard Model, for example, nobody would want to write every index, every generator and every tensor component explicitly. This would be a pointless and tedious task, and if not done with care and focus, it could even lead to unintentional errors. The toolkit is therefore intended to take a compact input and expand it according to the conventions and declarations encoded in the model.

Internally, the code will have to carry out a great deal of bookkeeping. In other words, it must be capable of expanding terms such as the covariant derivative according to the correct rules dictated by the symmetry groups that each Model possesses. It must be able to keep track of the order of the fields and their different indices. To do this, a sequence of small rewriting steps, each of which makes the expression slightly more explicit, is applied. One can visualise the structure as an hourglass: the user-facing interface allows several compact ways of declaring a Lagrangian, but these different inputs are progressively rewritten until they meet at the central object, the InteractionTerm. From this point on, the code can perform Wick contractions, canonicalisation and simplification of the resulting expression.

We want to proceed with an example in mind. Let's consider the quark kinetic term of QCD,
\begin{equation}
    i\,\bar q\,\gamma^\mu D_\mu q \;\supset\; g_s\,\bar q_{i}\,\gamma^\mu\,t^a_{ij}\,q_{j}\,G^a_\mu ,
    \label{eq:qcd-current}
\end{equation}
is, by the end of the pipeline, turned into the quark--gluon vertex $+i g_s\,\gamma^\mu\,t^a_{ij}$ in the conventions used here. The full gauge-fixed QCD Lagrangian,
\begin{equation}
\mathcal{L}_{\mathrm{QCD}}
= i\,\bar q\,\gamma^\mu D_\mu q
\;-\;\tfrac14 F^a_{\mu\nu} F^{a\,\mu\nu}
\;-\;\tfrac{1}{2\xi}\!\left(\partial^\mu A^a_\mu\right)^2
\;-\;\bar c^{\,a}\,\partial^\mu (D_\mu c)^a ,
\label{eq:qcd-full}
\end{equation}
is a good stress test because it exercises every branch of the compiler at once: a covariant matter kinetic term, a field strength, a gauge-fixing term and a ghost term.

% --------------------------------------------------------------
\section{Overview of the pipeline}
\label{sec:internals-overview}
% --------------------------------------------------------------

Let's start to look at the internals with a general overview, which is illustrated in
Figure~\ref{fig:pipeline-overview}. This shows the stages that lead from a declared \py{Model} to a Feynman-rule \py{Expression}. We can summarise it in four layers, see (Table~\ref{tab:internals-layers}).

\begin{figure}[htbp]
\centering
\begin{tikzpicture}[
    node distance=6mm,
    box/.style={draw, rounded corners, align=center, text width=8.4cm,
                minimum height=8mm, font=\small, inner sep=5pt},
    decl/.style={box, fill=blue!8},
    comp/.style={box, fill=orange!12},
    eng/.style={box, fill=green!12},
    post/.style={box, fill=purple!10},
    arr/.style={-{Latex[length=2.2mm]}, thick}
]

\node[decl]                      (model)     {\texttt{Model(...)}: declared fields, gauge groups, source Lagrangian};
\node[decl,  below=of model]     (declared)  {\texttt{DeclaredLagrangian.source\_terms} (still purely syntactic)};
\node[comp,  below=of declared]  (analyze)   {\texttt{\_analyze\_declared\_source\_term} $\to$ \texttt{AnalyzedSourceTerm} (routing)};
\node[comp,  below=of analyze]   (compile)   {\texttt{compiler.gauge}: expand \texttt{DC}, \texttt{FS}, gauge fixing, ghosts; re-lower};
\node[comp,  below=of compile]   (terms)     {tuple of \texttt{InteractionTerm} $\to$ \texttt{CompiledLagrangian}};
\node[eng,   below=of terms]     (legs)      {\texttt{feynman\_rule(...)}: build \texttt{ExternalLeg}s, \texttt{to\_vertex\_kwargs(...)}};
\node[eng,   below=of legs]      (contract)  {\texttt{contract\_to\_full\_expression}: sum over Wick contractions};
\node[post,  below=of contract]  (factor)    {\texttt{vertex\_factor}: overall $i$, optional momentum $\delta$, strip externals};
\node[post,  below=of factor]    (canon)     {idenso simplification + tensor canonicalisation};
\node[decl,  below=of canon]     (out)       {Symbolica \texttt{Expression}: the Feynman rule};

\foreach \a/\b in {model/declared, declared/analyze, analyze/compile,
                   compile/terms, terms/legs, legs/contract,
                   contract/factor, factor/canon, canon/out}
  \draw[arr] (\a) -- (\b);
\end{tikzpicture}
\caption{The pipeline from a declared \texttt{Model} to a Feynman-rule \texttt{Expression}. Blue = declaration layer (\texttt{feynpy}); orange = lowering/compiler layer (\texttt{feynpy.lowering}, \texttt{compiler.gauge} and helpers); green = contraction engine (\texttt{symbolic.vertex\_engine}); purple = post-processing (\texttt{symbolic.vertex\_postprocessing}, \texttt{symbolic.tensor\_canonicalization}).}
\label{fig:pipeline-overview}
\end{figure}

\begin{table}[htbp]
\centering
\begin{tabular}{@{}p{2.6cm}p{3.6cm}p{6.8cm}@{}}
\toprule
\textbf{Layer} & \textbf{Package(s)} & \textbf{Responsibility} \\
\midrule
Declaration &
\texttt{feynpy} (\texttt{metadata}, \texttt{declared}, \texttt{core}) &
Represents the physics as the user wrote it: declared fields, gauge groups, and source Lagrangian terms, still in a mostly syntactic form. \\
Lowering \& compilation &
\texttt{feynpy.lowering}, \texttt{lagrangian}, \texttt{compiler} &
Classifies and validates each source term, expands gauge structures such as covariant derivatives, field strengths, gauge-fixing terms and ghosts into local monomials, and collects the result into \texttt{InteractionTerm}s. \\
Contraction engine &
\texttt{symbolic.vertex\_engine} &
Builds external legs and vertex arguments, then sums over Wick contractions to assemble the full expression. \\
Post-processing &
\texttt{symbolic.vertex\_postprocessing}, \texttt{symbolic.tensor\_canonicalization} &
Applies the overall factor of $i$, optionally extracts the momentum-conserving $\delta$, strips external-leg factors, and simplifies and canonicalises the resulting tensor expression. \\
\bottomrule
\end{tabular}
\caption{The four layers of the pipeline, coloured as in Figure~\ref{fig:pipeline-overview}: declaration (blue), lowering/compilation (orange), contraction (green), and post-processing (purple).}
\label{tab:internals-layers}
\end{table}

The remaining sections walk the layers from the outside in: declaration (Section~\ref{sec:internals-declaration}), analysis and index resolution (Section~\ref{sec:internals-lowering}), gauge compilation (Section~\ref{sec:internals-compiler}), the \py{InteractionTerm} intermediate representation that ties everything together (Section~\ref{sec:internals-ir}), the Wick-contraction engine (Section~\ref{sec:internals-contraction}), and canonicalisation (Section~\ref{sec:internals-postprocessing}). Section~\ref{sec:internals-difficulty} then collects the hardest problems encountered while building the pipeline --- the places where the design had to be most deliberate.

\footnote{Each TikZ figure in this chapter has a plain-text counterpart kept as a Mermaid source diagram under \texttt{viz/mermaid/}}

% --------------------------------------------------------------
\section{The outer layer: declaring a model}
\label{sec:internals-declaration}
% --------------------------------------------------------------

\subsection{Inert metadata}

In the chapter \ref{ch:feynpy} we have seen all the building blocks of FEYNPY that will enable us to define the Model. These building blocks are the declarative \emph{metadata} stored in \texttt{feynpy.metadata}; these are frozen dataclasses that hold physics data and carry no algorithmic logic. Let's summarise them here:
\begin{itemize}
    \item \py{IndexType} pairs a Spenso representation with a ``kind'' string and a label prefix (more in Section~\ref{sec:cas-indextype});
    \item \py{Field} stores a name, spin, self-conjugacy, its tuple of \py{indices}, its Symbolica symbol and conjugate symbol, quantum numbers, and optional flavour-class data, inferring its \py{kind} (scalar/fermion/vector) and \py{statistics} from the spin;
    \item \py{Parameter} behaves algebraically like its symbol and may carry indices and explicit components;
    \item \py{GaugeRepresentation} and
    \item \py{GaugeGroup} record the symmetry data, whether a group is abelian, its coupling, gauge boson, ghost field, structure-constant builder and matter representations
\end{itemize}

The deliberate point is that metadata is inert: a faithful, machine-readable copy of what a physicist would write on paper, with all behaviour deferred downstream.

\subsection{The declarative DSL as wrapper objects}

Now that we have defined the building blocks of the theory, we can put them together into the Lagrangian itself. This is done with a small domain-specific language (DSL) defined in \texttt{feynpy.declared}:\footnote{A domain-specific language is a small language designed for a particular task. Here it is a restricted set of Python objects and functions that lets the user write Lagrangian terms in a notation close to standard physics notation.} \py{DC} (covariant derivative), \py{FS} (field strength), \py{Gamma}, \py{PartialD}, \py{Metric}, \py{T}, \py{StructureConstant}, \py{GaugeFixing}, \py{GhostLagrangian}. I want to emphasise that these constructors do not immediately build a Symbolica tensor expression. They build small immutable wrapper objects that will later be analysed and lowered into compiled interaction terms.

To understand how the code reads the Lagrangian let's consider as an example the equation \ref{eq:qcd-full}. When the user writes the equation \ref{eq:qcd-full} the first term looks like this
\begin{minted}{python}
I * Quark.bar * Gamma(mu) * DC(Quark, mu)
\end{minted}

Obviously the Lagrangian in equation \ref{eq:qcd-full} is composed of many terms, and so the first thing the code needs to do is to distinguish each monomial and separate its scalar coefficient from its field and tensor factors. More precisely the code assembles a \py{_DeclaredMonomial} which distinguishes between the coefficient and an ordered tuple composed of all the other factors. In our case its content is simply a scalar coefficient (\py{I}) and an ordered tuple that stores every factor --- here \py{_FieldFactor(Quark, conjugated=True)}, \py{GammaFactor(mu)}, \py{CovariantDerivativeFactor(Quark, mu)}. At this stage nothing is contracted, expanded or simplified. The order of the factors is preserved exactly because later passes need the original source structure. Later a process of analysis will determine what needs to be expanded.

A useful subtlety is that a few helpers are \emph{dual}. \py{Gamma(mu)} and \py{T(a)} are declarative placeholders: they say ``insert the gamma matrix'' or ``insert the generator'' and let the lowering pass infer the endpoints from the neighbouring fields. By contrast, \py{Gamma(i, j, mu)} and \py{T(a, i, j)} bypass that shorthand and return the fully explicit Spenso tensors \py{gamma(i, j, mu)} and \py{t(a, i, j)} immediately. This gives the user both a compact notation without specifying every index explicitly when the structure is obvious, and a fully explicit notation when it is not.

\subsection{\texttt{Model} and the two Lagrangian containers}

Once individual source terms have been written, the full source Lagrangian is stored in the container \py{DeclaredLagrangian}. This is a very simple object: essentially it is just an ordered tuple \py{source_terms}. In the most common case these source terms are \py{_DeclaredMonomial}s, but they can also be special declarations such as \py{GaugeFixing(...)} or \py{GhostLagrangian(...)}. So \py{DeclaredLagrangian} is nothing but the source-level Lagrangian, still written in the declarative language.

Everything is put together in the \py{Model} (in \texttt{feynpy.core}); this is the larger object that owns this declared Lagrangian together with the rest of the model data: \py{fields}, \py{parameters} and \py{gauge_groups}. Now the input is converted into a \py{DeclaredLagrangian}.

Actual compilation happens only when \py{model.lagrangian()} is called. At that point the declared source terms are turned into a \py{CompiledLagrangian}. This is where terms containing \py{DC(...)} or \py{FS(...)} are expanded using the metadata stored in the model, while terms that are already local are lowered directly. The result is no longer source syntax: it is a tuple of compiled \py{InteractionTerm}s. This compilation is lazy and cached, so it is only done the first time it is needed.

To summarise: \py{DeclaredLagrangian} is the source-level Lagrangian, \py{CompiledLagrangian} is the compiled one, and \py{Model} is the full object that stores the declared Lagrangian together with the extra information needed to validate and compile it.

% --------------------------------------------------------------
\section{Analysis and lowering: reading the Lagrangian and resolving indices}
\label{sec:internals-lowering}
% --------------------------------------------------------------

Between \py{DeclaredLagrangian} and \py{CompiledLagrangian} there is an intermediate step. The declared Lagrangian is the user-written Lagrangian, while the compiled Lagrangian is a tuple of local \py{InteractionTerm}s. Therefore we first need to make every implicit structure explicit, validate its index labels, and decide whether this term needs a model-dependent expansion. This process is what we call lowering.

For example, consider again the declared quark kinetic term
\begin{minted}{python}
I * Quark.bar * Gamma(mu) * DC(Quark, mu)
\end{minted}
Lowering checks that \py{mu} is used consistently as a Lorentz index and that the implied spinor and colour structure is well-formed, and classifies the term as a \py{covariant_core}. This still is not a final local interaction term, because \py{DC} must first be expanded using the model's gauge data. The compiler then performs that expansion, producing ordinary local monomials such as the kinetic piece and the quark--gauge-boson vertex, and those local monomials are what finally become compiled \py{InteractionTerm}s. By contrast, a source term that is already local, such as \py{lam * phi**4}, is lowered directly into its canonical interaction form with no further expansion.

So the lowering layer (\texttt{feynpy.lowering}) is the bridge between declarative syntax and the compiler. For each source term it answers three questions: what kind of term is this, are its indices well-formed, and (if the term is already local) what is its canonical interaction form?

\subsection{Classifying each source term}

In order to know what needs to be expanded, and how, by the compiler we want to analyse the terms. The single entry point is \py{_analyze_declared_source_term(...)}, which returns an \py{AnalyzedSourceTerm}. This object has the role of routing correctly every monomial through the compilation. If the term is already local, it already contains the final \py{InteractionTerm}; otherwise it stores the structured object that the compiler must expand. The compiler then consumes this analysed object, expanding it according to the gauge structure. The categories are summarised in Table~\ref{tab:analyzed-categories}.

\begin{table}[htbp]
\centering
\begin{tabular}{@{}p{4.0cm}p{5.0cm}p{3.2cm}@{}}
\toprule
\textbf{Category} & \textbf{Meaning} & \textbf{Example} \\
\midrule
\texttt{interaction} & a local \py{InteractionTerm}; no compilation & \py{lam * phi**4} \\
\texttt{covariant\_core} & canonical kinetic $\bar\psi\, i\gamma^\mu D_\mu \psi$ or $(D_\mu\Phi)^\dagger(D^\mu\Phi)$ & quark kinetic term \\
\texttt{generic\_covariant\_monomial} & any other monomial that still contains \py{DC} or another covariant operator needing expansion & current$\,\times\,$current \\
\texttt{gauge\_fixing} & a gauge-fixing declaration & \py{GaugeFixing(SU3, xi)} \\
\texttt{ghost} & a ghost declaration & \py{GhostLagrangian(SU3)} \\
\texttt{field\_strength\_monomial} & a monomial built from \py{FS(...)} factors & $-\tfrac14 F^a F^a$ \\
\bottomrule
\end{tabular}
\caption{The classification produced by \texttt{\_analyze\_declared\_source\_term}}
\label{tab:analyzed-categories}
\end{table}

So the analyser asks whether the term is already local; if not it proceeds with the classification. It checks if it is one of the two covariant kinetic shapes, afterwards if it is a generic covariant operator, which is a mix of covariant derivative, field strength and fields. There is also a dedicated classification for the gauge-fixing and ghost declarations (used in the declarative helper), and finally pure field-strength monomials.

\subsection{Resolving the indices}

As explained earlier, we do not want the user to have to write out every index in the Lagrangians, as many are implicit and obvious to a physicist. In this `lowering' stage, terms with implicit index structure are made explicit. The lowering process determines which symbolic labels belong to which field slots and which contractions the user intended to apply. A large amount of the design effort has been devoted to this stage, as it is precisely here that silent errors are most dangerous.


\paragraph{Monomial-wide label validation.} A first validation before lowering is done. The analyser builds one registry for every label \emph{name} appearing anywhere in the monomial (on field slots, on \py{PartialD}, \py{Gamma}, \py{T}, $f^{abc}$ and \py{FS}, and even on raw Spenso tensors in the coefficient) and requires that each name correspond to exactly one \py{IndexType}. Meaning that if the analyser encounters a label \py{f} as a flavour index it cannot appear on another field in the monomial as an adjoint index for example. So if the same name is used for two incompatible kinds the term is rejected with an explicit error.


\paragraph{Slot binding and inference.} Once each symbolic label has been assigned an index type (for example Lorentz, spinor, colour or flavour) the local monomial is resolved in a predefined order. Firstly, all labels that the user has entered directly into the fields are retained. Next, factors such as \py{Gamma} and \py{T}, which sit between two neighbouring fields, are resolved by identifying which field slots they connect to, to understand how to contract. Lowering replaces them by the corresponding explicit Spenso tensor, represented as a Symbolica expression, and multiplies that tensor into the coupling. Next, any other explicit labels --- such as those from declared factors or tensors already present in the coefficient --- are matched to the remaining free field slots of the same type. This matching must be one-to-one; if multiple assignments are possible, the lowering process generates an ambiguity error. After all explicit information has been used, the process proceeds to default assignment: it assigns new labels to obvious pairs, such as an adjacent $\bar\psi/\psi$, then to the sole remaining compatible pair, and finally introduces new dummy labels for any slots still left open. The result is that a compact source monomial is transformed into a fully explicit indexed interaction term, or rejected with a clear error if the notation is genuinely ambiguous.



\subsection{Remembering fermion pairings}

During lowering, the code records one further piece of information: how the fermion factors are paired into closed Dirac bilinears. This information is stored in the \py{closed_dirac_bilinears} metadata, a list of $(\bar\psi\text{-slot},\,\psi\text{-slot})$ pairs. The code determines these pairings automatically from the local spinor structure of the term.

This step is deliberately conservative. Lowering keeps this metadata only when the fermion slots can be partitioned unambiguously into disjoint, ordered closed bilinears. If no unique pairing is found, the monomial is rejected rather than guessed. When the pairing is unambiguous, the metadata preserves which \(\bar\psi\) is paired with which \(\psi\), so that the contraction engine can later apply the correct fermion-sign convention. Its importance is discussed in Section~\ref{sec:internals-difficulty}.

% --------------------------------------------------------------
\section{Compilation: turning gauge structure into ordinary interactions}
\label{sec:internals-compiler}
% --------------------------------------------------------------

Now the phase of full expansion of every operator of the Lagrangian is done. The compiler stage is the first stage that needs all the information stored in the model concerning the gauge symmetries. Through the lowering stage we analysed the terms and now with the compilation we turn the remaining unexpanded object into a local interaction term. The main dispatch lives in \texttt{src/compiler/gauge.py}: \py{_compile_declared_lagrangian_terms(model)} walks the analysed source terms and sends each one to a specialised handler (Figure~\ref{fig:compiler-dispatch}). Throughout this stage the conventions are
\begin{equation}
    D_\mu = \partial_\mu - i g\, A_\mu^a T_a,
    \qquad
    F_{\mu\nu} = \frac{i}{g}[D_\mu,D_\nu]
    = F^a_{\mu\nu} T_a ,
    \label{eq:conventions}
\end{equation}
with the adjoint basis fixed by
\begin{equation}
    [T_b,T_c] = i\, c^{a}{}_{bc}\, T_a,
    \qquad
    F^a_{\mu\nu}
    = \partial_\mu A^a_\nu - \partial_\nu A^a_\mu + g\, c^{a}{}_{bc}\, A^b_\mu A^c_\nu .
    \label{eq:conventions-components}
\end{equation}
For the usual \(\mathrm{SU}(N)\) basis one may simply read \(c^{a}{}_{bc}\) as the familiar \(f^{abc}\). Extra invariant tensors of the group, such as \(\epsilon\)-tensors or symmetric \(d\)-symbols, can certainly appear in local operators and coefficients, but they do not modify the definitions of \(D_\mu\) or \(F_{\mu\nu}\).

\begin{figure}[htbp]
\centering
\begin{tikzpicture}[
    node distance=5mm and 9mm,
    box/.style={draw, rounded corners, align=center, font=\footnotesize,
                inner sep=4pt, minimum height=7mm},
    comp/.style={box, fill=orange!12, text width=3.9cm},
    hand/.style={box, fill=orange!6, text width=4.6cm},
    ir/.style={box, fill=green!12, text width=2.6cm},
    arr/.style={-{Latex[length=2mm]}, thick}
]
\node[comp] (c1) {covariant\_core};
\node[comp, below=of c1] (c2) {generic\_covariant\_monomial};
\node[comp, below=of c2] (c3) {field\_strength\_monomial};
\node[comp, below=of c3] (c4) {gauge\_fixing};
\node[comp, below=of c4] (c5) {ghost};

\node[hand, right=of c1] (h1) {\texttt{\_compile\_declared\_covariant\_core}};
\node[hand, right=of c2] (h2) {\texttt{\_compile\_generic\_declared\_covariant\_monomial}};
\node[hand, right=of c3] (h3) {\texttt{compile\_field\_strength\_monomial}};
\node[hand, right=of c4] (h4) {\texttt{compile\_gauge\_fixing\_term}};
\node[hand, right=of c5] (h5) {\texttt{compile\_ghost\_term}};

\node[ir, right=of h3] (ir) {\texttt{InteractionTerm} tuple};

\foreach \c/\h in {c1/h1, c2/h2, c3/h3, c4/h4, c5/h5} \draw[arr] (\c) -- (\h);
\foreach \h in {h1, h2, h3, h4, h5} \draw[arr] (\h) -- (ir);
\end{tikzpicture}
\caption{Compiler dispatch. Each analysed category is handled by one routine, and every routine returns ordinary \texttt{InteractionTerm} objects. Structured operators are first expanded into local monomials and then lowered through the same path as hand-written interactions.}
\label{fig:compiler-dispatch}
\end{figure}

\subsection{Expansion and Interaction Term building}

The compiler essentially performs two operations. Firstly, it expands each structured gauge operator to yield simple local monomials. Secondly, only after carrying out the complete expansion does it feed these monomials back into the normal lowering process, generating a \py{InteractionTerm} (Section~\ref{sec:internals-lowering}).

This is a deliberate design choice: it prioritises generality and a single, uniform compilation pipeline over efficiency. For example dedicated routines could be written for $F^2$, $F^3$, $F^4$ and for each covariant structure separately, which would make it faster and more importantly closer to the canonical form we already know from textbooks. Here, however, there is no separate, hard-coded compiler for $F^2$, $F^3$, $F^4$ or for every possible power of a covariant derivative. The combinations are generated automatically via the standard distribution followed by a well-established lowering phase.

The cost is that the intermediate output is much more verbose than textbook notation. A non-abelian $F^2$ expands into nine local monomials before canonicalisation combines equivalent terms again (Section~\ref{sec:internals-postprocessing}). For higher operators this growth is rapid; for pure non-abelian $F^n$ it gives $3^n$ branches. Much of the effort has therefore gone into the canonicalisation process.

\subsection{Covariant derivatives}
\label{sec:internals-covd}

Let us examine the expansion of the covariant derivative in more detail. For a field $\phi$ transforming under several gauge groups, we define the general covariant derivative as

\begin{equation}
    D_\mu \phi
    =
    \partial_\mu \phi
    -
    i \sum_G g_G A^{a_G}_{G,\mu}\,
    \rho_G\!\left(T^{a_G}_G\right)\phi,
    \label{eq:covd-general}
\end{equation}

where the sum runs over all gauge groups $G$ under which $\phi$ transforms. Here, $T^{a_G}_G$ is a generator of $G$, while

\begin{equation}
    \rho_G : \mathfrak{g}_G \longrightarrow \operatorname{End}(V_G)
\end{equation}

is the representation map acting on the corresponding representation slot of $\phi$.

The two relevant cases are:

\begin{itemize}
    \item \textbf{Abelian group.} There is a single generator, represented by the charge $Q_G$. Therefore,
    \begin{equation}
        D_\mu \phi
        \supset
        -i g_G Q_G A^{(G)}_\mu \phi.
    \end{equation}

    \item \textbf{Non-abelian group.} The gauge field carries an adjoint index $a$, and the corresponding generator acts on the appropriate representation index of the field:
    \begin{equation}
        (D_\mu \phi)_i
        \supset
        -i g_G A^a_{G,\mu}
        \left(t^a_{G,R}\right)_{ij}\phi_j.
    \end{equation}
\end{itemize}

Let's take as an example a field that carries representation indices associated with several gauge groups, each generator acts only on its corresponding index. For example, for a field $\phi_{i\alpha}$ transforming in the tensor-product representation $R_1 \otimes R_2$,

\begin{equation}
    D_\mu \phi_{i\alpha}
    =
    \partial_\mu \phi_{i\alpha}
    -
    i g_1 A^a_\mu
    \left(t^a_{R_1}\right)_{ij}
    \phi_{j\alpha}
    -
    i g_2 B^b_\mu
    \left(t^b_{R_2}\right)_{\alpha\beta}
    \phi_{i\beta}.
\end{equation}

The routine \py{expand_cov_der} implements Eq.~\eqref{eq:covd-general} literally. It always emits one partial-derivative piece. It then resolves which gauge groups act on the field and emits one gauge piece for each of them; if the user did not name a group, it uses all compatible groups. For a conjugated occurrence $(D_\mu\phi)^\dagger$, the sign of the gauge piece flips.

\paragraph{More slots in the same representation.} The user might want to play with multiple slots in many representations that could be the same, in which case the generator may be able to act on more than one place. The code handles these cases. The default policy for this situation is conservative and rejects it as ambiguous.
 A representation can instead opt into \py{slot_policy='sum'}, in which case the compiler sums over the possible active slots.

\paragraph{Several gauge groups at once.} A standard example is the left-handed quark doublet $Q_L$, which transforms under $\mathrm{SU}(3)_c\times\mathrm{SU}(2)_L\times\mathrm{U}(1)_Y$. In that case
\begin{equation}
    D_\mu Q_L
    = \partial_\mu Q_L
    \;-\; i g_s\, G^a_\mu\, t^a_{(3)}\,Q_L
    \;-\; i g\, W^I_\mu\, t^I_{(2)}\,Q_L
    \;-\; i g'\, Y_{Q}\, B_\mu\, Q_L ,
    \label{eq:covd-multigroup}
\end{equation}
where we have $t^a_{(3)}$ as the colour generators, $t^I_{(2)}$ as the weak generators and $Y_Q$ as the hypercharge. So \py{DC(Q_L, mu)} expands to four pieces: the ordinary derivative plus one gauge current for each of the three groups. Naming one group, for example \py{DC(Q_L, mu, gauge_group=SU3C)}, keeps only that group's current together with the partial derivative.



\subsection{Field strengths}
\label{sec:internals-fs}

Let us examine the expansion of field strengths in more detail. As for the covariant derivative, the compiler expands \py{FS(...)} using the gauge-group data of the model so that it becomes a sum of local terms.

Using the convention

\begin{equation}
    D_\mu = \partial_\mu - i g A_\mu,
\end{equation}
with
\begin{equation}
    A_\mu = A_\mu^a T_a,
    \qquad
    [T_b,T_c] = i c^{a}{}_{bc} T_a,
\end{equation}
the Lie-algebra-valued field strength is
\begin{equation}
    F_{\mu\nu}
    =
    \partial_\mu A_\nu
    -
    \partial_\nu A_\mu
    -
    i g [A_\mu,A_\nu]
    =
    F^a_{\mu\nu}T_a,
\end{equation}
with components
\begin{equation}
    F^a_{\mu\nu}
    =
    \partial_\mu A^a_\nu
    -
    \partial_\nu A^a_\mu
    +
    g\,c^{a}{}_{bc}A^b_\mu A^c_\nu.
    \label{eq:fs-general}
\end{equation}


In the DSL, this is written as \py{FS(G, mu, nu, a)} for a non-abelian gauge group $G$, while for an abelian gauge group one writes \py{FS(G, mu, nu)}.

A single non-abelian \py{FS} factor is first rewritten as the sum
\begin{equation}
    F^a_{\mu\nu} = P_1 + P_2 + P_3,
\end{equation}
where
\begin{equation}
    P_1 = \partial_\mu A^a_\nu,
    \qquad
    P_2 = -\partial_\nu A^a_\mu,
    \qquad
    P_3 = g\,c^{a}{}_{bc}A^b_\mu A^c_\nu.
    \label{eq:fs-pieces}
\end{equation}
Note that for an abelian field strength, only $P_1$ and $P_2$ are present.

A monomial containing several field-strength factors is expanded by taking the full Cartesian product of all the different pieces.

To understand it more, let's consider the iconic example of the pure Yang--Mills Lagrangian
\begin{equation}
    \mathcal{L}_{\mathrm{YM}}
    =
    -\frac{1}{4}
    F^a_{\mu\nu}F^{a\,\mu\nu}.
\end{equation}
Each non-abelian field strength contributes three pieces, so their product generates
\begin{equation}
    3 \times 3 = 9
\end{equation}
raw branches before canonicalisation and collection. These branches belong to three sectors:
\begin{center}
\begin{tabular}{@{}lll@{}}
\toprule
\textbf{Piece combination} & \textbf{Count} & \textbf{Sector} \\
\midrule
$(P_1,P_2)\times(P_1,P_2)$
    & $2\times2=4$
    & quadratic kinetic sector \\
$(P_1,P_2)\times P_3$ and $P_3\times(P_1,P_2)$
    & $2+2=4$
    & cubic gauge interaction \\
$P_3\times P_3$
    & $1$
    & quartic gauge interaction \\
\bottomrule
\end{tabular}
\end{center}

Each resulting branch is then lowered in the same way as an ordinary local interaction.

Note that a product of $n$ non-abelian field strengths produces $3^n$ raw branches, whereas a product of $n$ abelian field strengths produces $2^n$ branches. More generally, a mixed product with $m$ non-abelian and $k$ abelian field strengths produces $3^m 2^k$ branches.



\subsection{Gauge fixing and ghosts}

There is a declarative helper for the gauge fixing term which can be accessed by \py{GaugeFixing(gauge_group, xi=...)}. This helper implements the linear covariant gauge term
\begin{equation}
    \mathcal{L}_{\mathrm{gf}}
    =
    -\frac{1}{2\xi}
    \left(\partial^\mu A^a_\mu\right)
    \left(\partial^\nu A^a_\nu\right),
\end{equation}
with the adjoint index omitted in the abelian case. It constructs the corresponding explicit local monomial, lowers it through the ordinary compilation pipeline, and tags the result as \py{gauge_fixing} so that the reporting and validation layers can recognise it. The user can naturally write the gauge fixing manually, in which case it is recognised and receives the same tag.

Note that the generic \py{GaugeFixing(...)} helper describes only linear covariant gauge fixing. More general $R_\xi$ gauges in spontaneously broken theories, for example, require additional Goldstone-dependent gauge-fixing terms, which must be written explicitly by the user.

There is also the ghost Lagrangian helper \py{GhostLagrangian(gauge_group)}. For a non-abelian gauge group, the Faddeev--Popov ghost Lagrangian can be written as
\begin{equation}
    \mathcal{L}_{\mathrm{gh}}
    =
    -\bar c^a\,
    \partial^\mu\left(D_\mu c\right)^a.
\end{equation}
Up to a total derivative, integration by parts gives
\begin{equation}
    \mathcal{L}_{\mathrm{gh}}
    \simeq
    \left(\partial^\mu\bar c^a\right)
    \left(D_\mu c\right)^a,
\end{equation}
where
\begin{equation}
    \left(D_\mu c\right)^a
    =
    \partial_\mu c^a
    +
    g\,c^{a}{}_{bc}A^b_\mu c^c.
\end{equation}
The ghost sector therefore becomes
\begin{equation}
    \mathcal{L}_{\mathrm{gh}}
    \simeq
    \left(\partial^\mu\bar c^a\right)
    \left(\partial_\mu c^a\right)
    +
    g\,c^{a}{}_{bc}
    \left(\partial^\mu\bar c^a\right)
    A^b_\mu c^c.
\end{equation}

The compiler turns the ghost sector into two ordinary local terms: the ghost kinetic term and the ghost gauge interaction term. Since the ghost field transforms in the adjoint representation, the code can use the same gauge-action and index-contraction machinery that it uses for other adjoint fields. What remains specific to ghosts is not the contraction structure, but their interpretation: they still carry Grassmann statistics, a fixed field ordering, and ghost number.

% --------------------------------------------------------------
\section{The pivot: the \texttt{InteractionTerm} intermediate representation}
\label{sec:internals-ir}
% --------------------------------------------------------------
To return to the analogy used at the start of the chapter, we are now at the waist of the hourglass. Every branch of the code converges and passes through this data type: \py{InteractionTerm} (in \texttt{feynpy.interactions}). Everything above this narrow point is declarative, everything below it is generic symbolic machinery.

\begin{minted}{python}
@dataclass(frozen=True)
class InteractionTerm:
    coupling: object                        # Symbolica Expression
    fields: tuple[FieldOccurrence, ...]     # bare field factors, in order
    derivatives: tuple[DerivativeAction, ...] = ()
    closed_dirac_bilinears: tuple[tuple[int, int], ...] = ()
    label: str = ""
    sector: str = ""
\end{minted}

The \py{coupling} is a Symbolica expression holding the numerical and tensor prefactor (couplings, generators, structure constants, metrics); \py{fields} is an ordered tuple of \py{FieldOccurrence} objects, each a field with its slot-resolved index labels and conjugation flag; \py{derivatives} records, as a \py{DerivativeAction}, which field each derivative hits and along which Lorentz index; and \py{closed_dirac_bilinears} carries the fermion-pairing provenance of Section~\ref{sec:internals-lowering}. A \py{CompiledLagrangian} is simply a tuple of such terms. An \py{InteractionTerm} is a purely static, index-labelled description of one monomial of the Lagrangian --- the direct code counterpart of the right-hand side
\[
    g_{\alpha_1\cdots\alpha_n}\,
    \partial\cdots\partial\,
    \phi^{\alpha_1}\cdots\phi^{\alpha_n}
\]
of Eq.~\eqref{eq:general-term}.

As we know the ordering of \py{fields} is relevant and since Symbolica multiplication is commutative, we use the ordered tuple which keeps record of fermion order.
The method \py{to_vertex_kwargs(external_legs)} performs the last model-aware step: it strips away all model-layer vocabulary and reduces one term plus a chosen set of external legs to flat, parallel lists --- species keys, momenta, field and leg roles, index labels and types, derivative indices and targets, and the bilinear metadata --- after which only Symbolica and Spenso objects remain.

\subsection*{Flavour expansion}

When the field carries a flavour index (Section~\ref{sec:feynpy-flavor}) by default the terms remain generic. For example $\lambda_{ij}\,\bar\ell_i \ell_j\,\phi$ will keep the lepton. By requesting \py{flavor_expand=True} FEYNPY enumerates every member combination, substitutes the corresponding \py{Parameter} component (e.g.\ $\lambda_{e\mu}$) producing one fully explicit \py{InteractionTerm} per combination in the same data structure.

\subsection*{Field transformations}
\label{sec:internals-transformations}

There is a second rewriting that lives at this layer, and it is the one that
produces the physical basis of the Standard Model. A
\py{FieldTransformation} (Section~\ref{sec:feynpy-transformations}) is applied
not to the declared source terms but to the compiled ones: each
\py{InteractionTerm} is visited, every \py{FieldOccurrence} matching a rule is
replaced by the corresponding expression, and the term is rebuilt.

Placing the stage here, rather than earlier, is a deliberate choice with two
consequences. First, \py{DC} and \py{FS} have already been expanded using the
gauge-group metadata of the original basis, which is the only basis in which
that metadata is meaningful --- $B_\mu$ has a hypercharge, the photon does
not. Second, because the transformation operates on an \py{InteractionTerm},
it inherits the structure that the term carries: derivatives recorded in
\py{derivatives} are distributed over a replacement product by the Leibniz
rule automatically, conjugated occurrences receive the conjugated
replacement, and the ordering that fixes fermion signs is preserved. A rule
that maps one field to a sum simply multiplies the number of terms, and
canonicalisation collects them again afterwards.

The transformation stage runs once, with all rules acting simultaneously, so
that a field introduced by one rule is not consumed by another in the same
pass. The result is an ordinary \py{CompiledLagrangian}, indistinguishable in
type from the one that entered, and vertex extraction proceeds from it
unchanged.

% -----------------------------------------------------
\section{The inner core: Wick contraction in the vertex engine}
\label{sec:internals-contraction}
% --------------------------------------------------------------

Now we arrive at the heart of the code, where the actual algorithm is implemented. We find the implementation of Wick contraction in
the vertex engine (\texttt{symbolic.vertex\_engine}). The only thing this engine cares about are fields, roles, index labels, momenta and a coupling tensor, and performs a combinatorial contraction that is the code-level realisation of the canonical-quantisation algorithm of Section~\ref{sec:feynrules-algorithm}.

\subsection{The permutation sum}

The core routine is \py{contract_to_full_expression}. Given an interaction with $n$ field slots and $n$ external legs, it sums over all $n!$ ways of matching legs to slots. Each permutation \py{perm} assigns leg \py{perm[i]} to slot \py{i}; the sum of the surviving contributions is the full Wick contraction.

For each permutation the code does the following things in order. First it remaps every open index in the coupling to the corresponding external-leg label, this ensures that the coupling will actually match the permuted fields. Then it checks whether the pairing is physically allowed, this is done using \py{factor_leg_compatible}. When it happens that a slot and a leg disagree in species, spin, statistics or conjugation, the permutation is discarded at once. Surviving fermionic terms pick up a Grassmann sign from the fermion part of the permutation, unless the closed-bilinear exception of Section~\ref{sec:internals-contraction-fermions} applies. Then derivatives are turned into momenta because of the Fourier rule $\partial_\mu \to -i p_\mu$ (all momenta incoming), so each \py{DerivativeAction} contributes a factor $(-i)\,p^\mu$ with $p$ the momentum of the matched leg (meaning that it will carry the corresponding label). Finally each leg contributes its external wavefunction ($U$ for bosons, $UF/\overline{UF}$ for fermions) and a species Kronecker delta, repeated indices are contracted with the Spenso metric, and the contribution is multiplied by the plane wave $\exp(-i\,\Sigma p\!\cdot\!x)$ before being added to the total.

Bosonic vertices carry no extra signs; fermionic ones do. The structure of the loop is:

\begin{minted}{python}
for perm in permutations(range(n)):
    term = coupling_with_open_indices_remapped_along(perm)
    if not all(factor_leg_compatible(i, perm[i], ...) for i in range(n)):
        continue                       # prune impossible pairing
    if statistics == "fermion" and not use_closed_dirac_bilinear_signs:
        term *= fermion_sign(perm, fermion_slots)
    for mu, tgt in zip(derivative_indices, derivative_targets):
        term *= (-I) * pcomp(ps[perm[tgt]], mu)
    for i, j in enumerate(perm):
        term *= external_factor(role[i], alphas[i], betas[j], ps[j], ...)
    term *= plane_wave(sum_of_matched_momenta, x)
    total += term
\end{minted}

Note that nothing here depends on a particular index kind. An index is only a \py{(kind, label)} pair.

\subsection{Fermion signs and the closed-bilinear exception}
\label{sec:internals-contraction-fermions}

Getting fermion signs right was one of the trickier parts of the engine. The pipeline records the order of the fields from the very start, so that even this final layer still knows which sign to apply.

The default rule is simple: use the flat Grassmann sign of the fermion permutation. This works in general, but it fails for a product of closed Dirac bilinears. A concrete example is the identical-fermion current--current operator
\begin{equation}
    (\bar\psi\,\gamma^\mu\psi)\,(\bar\psi\,\gamma_\mu\psi).
\end{equation}
This operator has two Wick contractions: a \emph{direct} one, where each current stays intact, and a \emph{crossed} one, where the two currents exchange a fermion. The flat Grassmann rule predicts a relative sign of $\text{direct} - \text{crossed}$ between the two, but the correct sign, following the FeynRules convention, is $\text{direct} + \text{crossed}$.


The fix uses the fermion-pairing information recorded during lowering (Section~\ref{sec:internals-lowering}). If that information shows that every Dirac fermion slot belongs to exactly one closed bilinear, the engine drops the flat permutation sign and treats every pairing with a plus. In every other case (for example a generic open multi-fermion tensor) the flat Grassmann sign still applies.

\subsection{Output policy and universal factors}

There is a function, namely \py{vertex_factor}, that has the role of running the contraction and then applies a few presentation choices. By default it amputates the external wavefunctions $U$, $UF$ and $\overline{UF}$, so the user sees the bare vertex. Optionally, the user can choose to write \py{include_delta=True}, which gives the momentum-conservation factor $(2\pi)^d\,\delta(\Sigma p)$. In every case the result is finally multiplied by the overall $+i$ that every Feynman rule carries. These choices only decide how the answer is shown.

% --------------------------------------------------------------
\section{Making the answer canonical}
\label{sec:internals-postprocessing}
% --------------------------------------------------------------

The raw contracted expression is correct, but not unique. The same vertex can appear with different dummy-index names, different orderings of antisymmetric tensors, or still expanded into many field-strength pieces. In the next paragraph some choices of the postprocessing of the final output are presented.

The first pass, in \texttt{symbolic.vertex\_postprocessing}, does the physics-facing clean-up. As the output comes out of the contraction engine it has the full range of Kronecker deltas of the permutation. The first step is therefore to contract species Kronecker deltas, collapse repeated bispinor indices with the Spenso metric, and, on request, simplify closed gamma chains with idenso (Section~\ref{sec:cas-idenso}).

The second pass, in \texttt{symbolic.tensor\_canonicalization}, puts the surviving expression into a unique tensor form using Symbolica's \py{canonize_tensors} (Section~\ref{sec:cas-symbolica}). That routine only understands symmetry flags on plain symbols, not on Spenso tensor heads. So the code temporarily replaces every Spenso tensor --- metrics, generators, structure constants, $\gamma$ matrices, and the project's custom invariants --- by a plain symbol with the right symmetry, runs the canonicalisation, and swaps the Spenso tensors back. Dummy indices of different kinds (Lorentz, colour, spinor, weak) are kept in separate groups, so a Lorentz dummy is never confused with a colour one.

On top of this, \py{canonize_full} adds a few physics reductions used when checking identities such as gauge variation and BRST (Sections~\ref{sec:feynpy-gauge-variation} and~\ref{sec:feynpy-brst}): it makes flat derivatives commute, reduces products of structure constants with the Jacobi identity, and drops Yang--Mills terms that vanish by the antisymmetry of $f^{abc}$.

Together these passes turn the over-expanded compiler output into a compact, comparable Feynman rule.

% --------------------------------------------------------------
\section{Where the difficulty lived}
\label{sec:internals-difficulty}
% --------------------------------------------------------------

Building the pipeline was less about one clever algorithm than about a few recurring hard problems. These are the places where most of the design effort went, and they shaped the architecture.

\paragraph{Index bookkeeping.} The hardest bugs came from indices, not from algebra. Early code keyed labels only by index kind, which fails as soon as a field carries two slots of the same kind --- two Lorentz indices, or two colour indices. The fix is now layered: labels are packed by slot, each label name is bound to one index type across the whole monomial, ambiguous assignments are rejected rather than guessed, and gauge representations expose an explicit \py{slot_policy}. The guiding rule is to fail early: an ambiguous index assignment is an error at declaration time, never a silent wrong contraction.

\paragraph{Fermion signs.} Matching the FeynRules convention for products of closed Dirac bilinears, without breaking generic multi-fermion terms, required carrying the bilinear pairing from lowering all the way into the contraction engine (Section~\ref{sec:internals-contraction-fermions}). Getting that boundary right --- all pluses only when every fermion sits in a closed pair --- took careful checks on benchmark operators and is now a fixed convention.



\paragraph{Making expansion generic.} Replacing a special-case $F^2$ handler with a generic \py{FS} expansion --- and likewise for covariant derivatives --- unlocked $F^3$, $F^4$ and mixed-group operators for free. The cost was many more raw local monomials, and a performance problem in the contraction engine: the old compatibility check compared only broad field roles and therefore tried many impossible pairings. The species-aware pruning of Section~\ref{sec:internals-contraction} fixed that without changing any resulting formula.

\paragraph{Keeping the layers honest.} Finally, much of the work went into not letting the layers leak into each other: declarative syntax must never reach the contraction engine, the engine must never know what a covariant derivative is, and comparison against external references must live outside the engine. One job per layer, and one intermediate representation at the waist, is what made it possible to validate the packaged Standard Model against the reference export vertex by vertex.

% --------------------------------------------------------------
\section{Summary}
\label{sec:internals-summary}
% --------------------------------------------------------------

The pipeline is a chain of small, inspectable rewrites. A \py{Model} holds the source Lagrangian and the inert data of the theory. Lowering classifies each term, checks its indices, and turns local terms into \py{InteractionTerm}s. The compiler expands gauge structure --- covariant derivatives, field strengths, gauge fixing and ghosts --- into ordinary local monomials and lowers them through the same path. Every branch therefore meets at the same intermediate representation. The symbolic engine then contracts it by a pruned permutation sum, applies momenta, external factors, fermion signs and the overall $+i$, and canonicalisation makes the answer unique. The hard parts --- indices, fermion signs, conventions, generic expansion and layer boundaries --- are exactly where the design had to be most careful, and they are what the validation in the next chapter rests on.
